\section{Discussion}
\label{sec:discussion}

\subsection{Interpretation}

The coverage result is structural, following from what each evidence class can
physically observe rather than from a threshold setting, and it was obtained
with calibration held fixed across all scenarios. The implication for the
literature is narrow but practical. An accuracy reported against
localisation-visible attacks is an accuracy over that subset of the attack
space. EMBN-MANET's threat model is confined to false-location behaviour by
design, so its figure is correct within its scope and we do not dispute it.
What we would ask of any such report is one additional line: which attack
classes the evidence can observe at all.

The isolation result is the paper's strongest. Over 30 seeds,
confirmation-gated isolation withdrew 3.9\% of honest-beacon service on an
uncontested network (5.0\% with evidence-aware release) against 74.8\% for
the strongest immediate policy, with every paired interval excluding zero; the
ordering held under attack, under shadowing and when the schedules were
enforced. With evidence-aware release, its cost is the time an attacker
remains active before confirmation: a median of 15\,s for attacks with a localisation
signature and 27.5\,s for selective forwarding, rising to medians of 125--155\,s for
forwarding attacks under shadowing. Whether that trade is favourable depends on
the prior probability of attack and on what a wrongly removed beacon costs the
network. Where attack is rare and beacon coverage sparse, gating is preferable,
since $\Phi$ accumulates continuously while attacker exposure is bounded by the
confirmation delay. That bound holds only if release cannot be triggered by
missing evidence: under plain E a forwarding attacker is re-admitted once its
isolation silences its behavioural branch. Release is as consequential a
decision as isolation, and in our data the safe release rule made isolation on
confirmation close to permanent. Where attack is frequent and beacon redundancy high,
immediate isolation may be preferable. We present the trade rather than a
winner, and suggest that $\Phi$ or a comparable measure be reported routinely:
a mechanism's false-positive rate says nothing about how long its errors
persist.

\subsection{Limits of the evidence}

Whether fusion helped was decided by the attack family, not by the
combination rule. In the primary sweep the branches' sensitivity ratio was 3.63, beyond the
bound of 2.414, and Stouffer combination confirmed fewer attackers than a
disjunctive rule at the two intensities nearest the behavioural threshold.
Halving the drop rate relative to displacement brought the ratio to 1.87,
inside the bound, and fusion then never confirmed fewer and confirmed one more.
The detector was the same in both sweeps; only the mix of displacement and
drop rate changed. Fusion contributes against attackers whose two
misbehaviours are comparably visible to the two branches, and against
lopsided ones a disjunctive rule does as well or better. The baseline
comparison is the same result at the level of whole detectors: in the small
dense network, two calibrated single-evidence detectors run in parallel
matched the fused pipeline on every attack evaluated and were faster under
shadowing. In the sparse and 50-node networks the behaviour-only detector's
false isolation rose to 11--24\% of honest-beacon service while the
pipeline's stayed at 0--3\%, plausibly because of its confidence weighting
and peer-relative scoring;
there the benefit of fusion appears in $\Phi$, not in detection. Branch sensitivity to
a given attack family can be measured before fusion is applied, and no system
we surveyed reports checking it. An adaptive adversary could exploit this by
concentrating its misbehaviour on one axis, which is a further reason to
keep a single-branch path to confirmation alongside the fused statistic.

The aggregation stage changed no verdict anywhere. It belongs in the framework
as an unevaluated robustness extension, not as a validated component: a system
using only direct per-observer trust would reproduce every reported finding.
It remains because our campaign barely samples the conditions it exists to
handle, namely observers who disagree, direct trust that is noisy, and
recommenders who lie.

Neither isolation policy produced a defensible throughput benefit. This may be
a property of a dense 25-node topology in which AODV finds alternative routes
readily; enforcing isolation in the sparse topology is the obvious next
experiment.
Multipath recovery schemes~\cite{sirajuddin2021tbsmr} and isolation studies in
related ad hoc settings~\cite{kavitha2019india,buksh2024trust} suggest the
effect may be topology-dependent. On the evidence we have, the claim that
isolation restores network performance is unsupported. Finally,
throughput-based evaluation of A1 measures nothing, since A1 has no
packet-level footprint; the appropriate metric for that attack class is
localisation error.

\section{Threats to Validity}
\label{sec:threats}

\textbf{Construct validity.} Behavioural monitoring is node-global rather than
per-observer, so indirect trust operates on a coarser behavioural signal than
the model specifies. This weakens the aggregation ablation: the null may
reflect the implementation rather than the design. $\Phi$ for the 30-seed
campaign is computed by applying each policy offline to the same verdict
stream, which assumes that an isolation does not change later verdicts. For
the never-releasing policies the first isolation time is exact and only later
isolations in the same run can be affected; for E the release times can be,
and Section~\ref{sec:release} shows that for forwarding attackers they are.
Enforcing the replayed schedules on ten clean seeds and rebuilding them from
the verdicts issued under isolation moved neither policy's $\Phi$ by an amount
distinguishable from zero and preserved the D$_1$--E gap in every seed
(Section~\ref{sec:phi}). Iterating the loop for three rounds left E$'$ at a
fixed point on every clean seed and never released an attacker, while E
re-admitted the A2 attacker in every seed. Iteration approximates an online
detector but is not one, and it covers E and E$'$ on two scenarios. E$'$ was evaluated from verdicts
issued before isolation; its schedule is exact for beacons confirmed on
behaviour, which it never releases, and inherits the replay assumption for
beacons confirmed on localisation alone. E$'$ was designed after the
enforced runs were examined, so its evaluation is not blind to that
finding.

\textbf{Internal validity.} The recovery baseline window is short relative to
throughput variance, which is why we exclude throughput debt. Detection
thresholds were calibrated on held-out clean runs and fixed, which avoids
tuning to the test scenarios but means no scenario received its optimal
operating point.

\textbf{External validity.} The primary channel has no shadowing, multipath or
obstruction model; the shadowing campaign shows that detection and the
false-isolation result survive 4--8\,dB, at a substantial cost in latency for
forwarding attacks. The campaign is otherwise narrow in five ways. Attackers
are non-colluding and non-adaptive. Mobility follows a synthetic model rather
than a trace. The topologies are a 25-node network in a dense and a sparse
field and a 50-node network at the dense field's density; the sparse and
50-node campaigns measured detection and false isolation but not the
throughput effect of enforced isolation, so the throughput nulls rest on the
small dense field alone. Larger networks were limited by memory: a 50-node
run needs about 3\,GB, and a 100-node configuration did not fit the hardware
available. Three monitored beacons per run, six at 50 nodes, make rate-based
metrics coarse. And
only the hybrid attack's intensity is swept. We do not claim the results
generalise beyond this configuration.

\textbf{Comparative validity.} The two literature baselines re-implement
the evidence class and decision rule of Liu--Ning--Du alert counting and of
watchdog/CONFIDANT reputation; they are not complete reimplementations of
those systems, and they use our simulator's node-global behavioural counters
and our residual. Their thresholds were set on calibration seeds by rules we
chose, and a different but equally defensible rule would move their operating
points. Policies D, D$_w$ and D$_1$ are our own constructions, reasonable
stand-ins for isolate-on-suspicion behaviour. The comparison against
EMBN-MANET itself remains analytical, resting on a reading of its evidence
class and threat model.

\textbf{Statistical conclusion validity.} Coverage, latency, false positives
and $\Phi$ rest on 30 seeds with held-out calibration and paired bootstrap
intervals. The throughput, recovery, localisation-error and ablation results
rest on the five seeds of the original campaign; no difference among the
continuous throughput measures is claimed. The sparse and 50-node campaigns use ten seeds each, so their paired
intervals are wide and two of the 50-node comparisons reach zero. The sweeps
use ten seeds per level,
so the one-seed fusion gain in the second pairing is reported as consistent
with the condition, not as a measured benefit.

\section{Conclusion}
\label{sec:conclusion}

We evaluated multi-evidence verification and confirmation-gated isolation of
malicious beacon nodes in a mobile ad hoc network across 780 ns-3 simulation
runs.

For the two attack classes evaluated, single-evidence detection is blind to
attacks outside its evidence class, and collecting both behavioural and
localisation evidence removes the observed blind spots. Confirmation-gated
isolation reduced cumulative false isolation of honest beacons to 3.9\% of
available service against 74.8\% for the strongest immediate-isolation policy
over 30 seeds, at a median cost of 15--27.5\,s of attacker activity after
onset, and the gap held when the schedules were enforced. Releasing a beacon
when its verdict clears re-admitted forwarding attackers, whose behavioural
evidence disappears while they are isolated; requiring that evidence to return
before release prevented every re-admission and raised false isolation only to
5.0\%. Both results held under 4 and 8\,dB shadowing, where the price moved to
latency: forwarding attacks took a median of 125--155\,s to confirm. Every weighted
Stouffer detector has a fusion-only region, but an attack family scaled in
intensity reaches it only when the branches' detection thresholds for that
family lie within a factor of $1+\sqrt2$, and across two sweeps that condition
predicted whether fusion diluted ($r = 3.63$) or matched and marginally
exceeded ($r = 1.87$) a disjunctive rule. Consistent with that condition, a
calibrated disjunction of localisation-only and behaviour-only literature
baselines matched the fused pipeline's coverage and false isolation and
isolated attackers as fast or faster in the small dense network. In the
sparse network the disjunction detected the forwarding attacker more often
but withdrew 24\% of honest-beacon service, against none for the pipeline,
and in a 50-node network with two attackers it withdrew 11\% against 3\%
while the pipeline confirmed every attacker. The trust aggregation layer changed
no verdict, and neither isolation policy produced a throughput difference
exceeding its confidence interval.

We report these outcomes together because the negative ones bound the positive
ones. In a literature where reported accuracies cluster above 95\%, the more
informative questions are which attacks a detector can see and how long its
mistakes last.

\subsection{Future work}
\label{sec:future}

\textbf{Complete baselines.} Full implementations of EMBN-MANET and of a
behavioural trust-routing scheme, including their response mechanisms, on the
same topology and traces.

\textbf{Scale and topology.} Networks of 100 nodes and more, enforced
isolation in sparse topologies where it can affect delivery, varied mobility
speeds, and colluding or adaptive attackers.

\textbf{Fusion-rule comparison.} Bounded rules (arithmetic and geometric mean),
calibrated statistical rules (Stouffer, Fisher, Dempster--Shafer) and learned
rules (regularised logistic regression, gradient-boosted trees) compared on
identical oriented evidence, in configurations chosen on both sides of the
bound~\eqref{eq:ratio}, with folds partitioned by seed.

\textbf{Aggregation in its intended regime.} Noisy observers, conflicting
observers and malicious recommenders, which is where indirect and historical
trust are designed to matter.

\textbf{Post-isolation localisation error.} Splitting localisation error at the
isolation timestamp would give direct evidence that isolation restores a
degraded service, which throughput cannot show for A1.

\textbf{Online closed-loop evaluation.} The policies here were evaluated by
replay, one-step enforcement and a three-round iterated loop. Running the detector inside the
simulator, so that each isolation and release feeds back into later
verdicts, would test E$'$ and the replay assumption directly, and would allow
probationary release, in which a released beacon is re-isolated on its first
malicious window.

\textbf{Faster behavioural confirmation under shadowing.} The latency cost of
forwarding attacks under shadowing comes from the behavioural branch's slow
confidence saturation and widened peer distribution; adaptive half-weights or
a second behavioural observable could reduce it.

\section*{Acknowledgements}

The authors thank the developers and maintainers of the ns-3 network
simulator~\cite{ns3}, whose simulator and AODV, IEEE 802.11 and propagation
models this study builds on, and the developers of the open-source Python
scientific stack used for the analysis.

\section*{Declarations}

\textbf{Funding.} No funding was received for this work.

\textbf{Competing interests.} The authors declare no competing interests.

\textbf{Author contributions.} A.V.\ designed and implemented the simulator,
verification pipeline and experiments, and drafted the manuscript. D.V.C.\
supervised the work, contributed to the experimental design and revised the
manuscript. Both authors read and approved the final version.

\textbf{Data and code availability.} Simulation sources, the analysis
pipeline, the scripts that run every campaign and regenerate every data
figure, and the result tables, per-window verdicts and isolation schedules
for every campaign are available at
\url{https://github.com/kraghu356/manet-beacon-trust}. The raw per-run
simulator traces are available from the corresponding author
on reasonable request. No third-party dataset was used: every data point in
this paper was generated by the authors' simulations.

\textbf{Ethics approval.} Not applicable. This study involves neither human
participants nor animals.
